where $U_k = S_x (I_k \otimes C_x)$ and $U_m = S_y (I_m \otimes C_y)$
are independent 1D evolution operators on the $x$- and $y$-cycles.
The factorisation isolates the spectral constraints along each axis,
so the root-of-unity conditions that govern full-state revivals can
be evaluated dimension-by-dimension.

A torus revival requires the sum of the two axis phases to be zero. When both axes revive 
at a common step, the sum vanishes at the LCM of the two axis periods; but two nonzero 
phases of opposite sign can also sum to zero at a step where neither axis has revived on its 
own, which gives a torus revival at a period shorter than the LCM of the two axis periods.

\emph{Note on exact versus projective revivals.}

In all cases a revival requires both the position and spin state of the walker to
return to its initial values.  In this work we distinguish between two classes of revivals:
a \emph{projective} revival is one in which the state returns up to an
unobservable global phase ($U^N = e^{i\phi} I$).  an \emph{exact}
or full-state revival is one where the state returns with no phase ($U^N = I$). 
Here we enumerate only exact revivals. 
